% Part: first-order-logic
% Chapter: introduction
% Section: semantic-notions

\documentclass[../../../include/open-logic-section]{subfiles}

\begin{document}

\olfileid{fol}{int}{sem}

\olsection{অর্থগত ধারণা}

উপরে বলেছি, $\Sat{M}{!A}[s]$ সিদ্ধ হয় কি না বিবেচনার সময় সুবিধার জন্য
$s$-কে সব !!{variable}s-এর মান দিতে দিই; কিন্তু~$!A$-তে থাকা
!!{variable}s-কে সেটি যে মানগুলি দেয়, শুধু সেগুলিই ব্যবহৃত হয়। বস্তুত,
$!A$-তে \emph{মুক্ত} চলরাশিগুলির মানই কেবল প্রাসঙ্গিক। অবশ্যই, আমরা
সতর্ক বলে এই তথ্যটি প্রমাণ করব। !!{sentence}s-এ কোনো মুক্ত চলরাশি নেই;
তাই কোনো !!a{structure}-এ সেগুলি পরিতৃপ্ত হয় কি না, তার ক্ষেত্রে~$s$-এর
কোনো প্রভাবই নেই। অতএব $!A$ যখন !!a{sentence}, তখন “সব~$s$-এর জন্য
$\Sat{M}{!A}[s]$”-কে $\Sat{M}{!A}$ বলে সংজ্ঞায়িত করতে পারি; ঘটনাক্রমে
এটি সিদ্ধ হয় যদি এবং কেবল যদি অন্তত একটি~$s$-এর জন্য
$\Sat{M}{!A}[s]$ সিদ্ধ হয়। !!{formula}s-এর পরিতৃপ্তির কার্যকর সংজ্ঞা
পেতে !!{variable} আরোপ প্রবর্তন করতে হয়; কিন্তু !!{sentence}s-এর
পরিতৃপ্তি !!{variable} আরোপ থেকে স্বাধীন।

“$\Sat{M}{!A}$”-এর সংজ্ঞা পাওয়ার পর প্রথম-ক্রমের যুক্তিবিদ্যার
!!{sentence}s-এর ক্ষেত্রে “ক্ষেত্র” ও “এর মধ্যে সত্য”-এর অর্থ জানা যায়।
!!{sentence}s-এর জন্য $\Sat{M}{!A}$-র সংজ্ঞার ভিত্তিতে বৈধতা, অনুসিদ্ধান্ত
ও পরিতৃপ্তিযোগ্যতার মৌলিক অর্থগত ধারণাগুলি সংজ্ঞায়িত করা যায়। কোনো
বাক্য বৈধ, $\Entails !A$, যদি প্রতিটি !!{structure} তাকে পরিতৃপ্ত করে।
!!{sentence}s-এর কোনো সেট থেকে সেটি অনুসৃত হয়, $\Gamma \Entails !A$,
যদি~$\Gamma$-র সব !!{sentence}s-কে পরিতৃপ্ত করা প্রতিটি !!{structure}
$!A$-কেও পরিতৃপ্ত করে। আর !!{sentence}s-এর কোনো সেট পরিতৃপ্তিযোগ্য, যদি
কোনো !!{structure} তার সব !!{sentence}s-কে একই সঙ্গে পরিতৃপ্ত করে।

!!{formula}s আরোহীভাবে সংজ্ঞায়িত এবং পরিতৃপ্তিও !!{formula}s-এর গঠনের
উপর আরোহের সাহায্যে সংজ্ঞায়িত বলে, আমাদের অর্থতত্ত্বের ধর্মগুলি প্রমাণ
করতে ও সংজ্ঞায়িত অর্থগত ধারণাগুলির মধ্যে সম্পর্ক স্থাপন করতে আরোহ
ব্যবহার করা যায়। আমরা এসব ধর্মের কয়েকটি একত্র করে প্রমাণ করব—কিছুটা
তাদের প্রত্যেকটির নিজস্ব গুরুত্বের জন্য, তবে প্রধানত অর্থতত্ত্ব ও
!!{derivation} পদ্ধতির সম্পর্ক অনুসন্ধান করতে গেলে তাদের অনেকগুলি কাজে
লাগবে বলে। তা করতে সংকেতবিন্যাসের এমন কয়েকটি ধারণা ও ক্রিয়াকেও
(নির্ভুলভাবে, অর্থাৎ আরোহের সাহায্যে) সংজ্ঞায়িত করতে হবে, যেগুলির উল্লেখ
এখনও করিনি।

\end{document}
